\documentclass[12pt,a4paper]{article}
\usepackage[utf8]{vietnam}
\usepackage{geometry}
\geometry{left=2.2cm,right=2.0cm,top=2.2cm,bottom=2.2cm}
\usepackage{setspace}
\setstretch{1.18}
\usepackage{microtype}
\usepackage{xcolor}
\definecolor{mainblue}{HTML}{1D4ED8}
\definecolor{softblue}{HTML}{EFF6FF}
\definecolor{softgreen}{HTML}{ECFDF5}
\definecolor{softorange}{HTML}{FFF7ED}
\definecolor{softred}{HTML}{FEF2F2}
\definecolor{darkgray}{HTML}{111827}
\definecolor{midgray}{HTML}{4B5563}
\usepackage{xurl}
\usepackage{amsmath,amssymb}
\usepackage{booktabs}
\usepackage{longtable}
\usepackage{tabularx}
\usepackage{array}
\newcolumntype{L}[1]{>{\raggedright\arraybackslash}p{#1}}
\usepackage{multirow}
\usepackage{graphicx}
\usepackage{enumitem}
\setlist[itemize]{leftmargin=1.2cm,itemsep=2pt,topsep=2pt}
\setlist[enumerate]{leftmargin=1.2cm,itemsep=2pt,topsep=2pt}
\renewcommand{\arraystretch}{1.18}
\usepackage{tikz}
\usepackage{pgfplots}
\pgfplotsset{compat=1.18}
\usetikzlibrary{arrows.meta,positioning,shapes.geometric,shapes.misc,fit,calc}
\usepackage[most]{tcolorbox}
\usepackage{fvextra}
\DefineVerbatimEnvironment{Code}{Verbatim}{fontsize=\small,breaklines=true,breakanywhere=true,frame=single,framesep=2mm,rulecolor=\color{gray!55}}
\newtcolorbox{infobox}[1][]{breakable, colback=softblue, colframe=mainblue!50!black, boxrule=0.5pt, arc=2mm, left=1mm, right=1mm, top=1mm, bottom=1mm, #1}
\newtcolorbox{warnbox}[1][]{breakable, colback=softorange, colframe=orange!70!black, boxrule=0.5pt, arc=2mm, left=1mm, right=1mm, top=1mm, bottom=1mm, #1}
\newtcolorbox{redbox}[1][]{breakable, colback=softred, colframe=red!70!black, boxrule=0.5pt, arc=2mm, left=1mm, right=1mm, top=1mm, bottom=1mm, #1}
\newtcolorbox{notebox}[1][]{breakable, colback=softblue, colframe=mainblue!50!black, #1, boxrule=0.5pt, arc=2mm, left=1mm, right=1mm, top=1mm, bottom=1mm}
\newtcolorbox{keybox}[1][]{breakable, colback=softgreen, colframe=green!45!black, #1, boxrule=0.5pt, arc=2mm, left=1mm, right=1mm, top=1mm, bottom=1mm}
\newtcolorbox{riskbox}[1][]{breakable, colback=softred, colframe=red!55!black, #1, boxrule=0.5pt, arc=2mm, left=1mm, right=1mm, top=1mm, bottom=1mm}
\DeclareRobustCommand{\code}[1]{\begingroup\urlstyle{tt}\url{#1}\endgroup}
\newcommand{\term}[1]{\textbf{#1}}
\DeclareRobustCommand{\file}[1]{\code{#1}}
\newcommand{\ra}{$\rightarrow$}
\sloppy
\usepackage{titlesec}
\titlespacing*{\section}{0pt}{1.2em}{0.55em}
\titlespacing*{\subsection}{0pt}{0.9em}{0.35em}
\titlespacing*{\subsubsection}{0pt}{0.7em}{0.25em}

\begin{document}

\begin{titlepage}
    \centering
    \vspace*{1.6cm}
    {\Huge\bfseries Kịch bản demo Phân hệ 2\\[0.35cm]}
    \vspace{0.7cm}
    {\Large Bảo mật dữ liệu y tế trên Oracle — chia việc cho 5 người\par}
    \vspace{1.2cm}
    \begin{infobox}
    \textbf{Phạm vi:} Kịch bản trình diễn phủ \textbf{TC\#1} (ánh xạ tài khoản), \textbf{YC1} (RBAC + VPD),
    \textbf{YC2} (OLS), \textbf{YC3} (Audit), \textbf{YC4} (Backup/Recovery) và \textbf{phần mở rộng mã hóa}
    (NNE đường truyền, TDE at-rest, masking, đổi mật khẩu). Bám theo \file{DEMO_SCRIPT.md} và \file{TALKING_POINTS.md}.
    \end{infobox}
    \vfill
    {\large Môn CSC12001 — An toàn \& Bảo mật Hệ thống Thông tin.\par}
    \vspace{0.6cm}
    {\large \today\par}
\end{titlepage}

\pagenumbering{roman}
\tableofcontents
\clearpage
\pagenumbering{arabic}

\section{Mục tiêu \& thông điệp xuyên suốt}

Hệ thống là ứng dụng WinForms (.NET) kết nối Oracle. Đồ án không để ứng dụng tự quyết định quyền
truy cập dữ liệu y tế; toàn bộ RBAC, VPD, OLS, Audit và trigger được thực thi tại Oracle Database
Engine. Mỗi người trình bày một phần, nhưng tất cả nên nhắc lại cùng một thông điệp.

\begin{keybox}[title={Thông điệp xuyên suốt (mọi người nhắc lại)}]
Ranh giới bảo mật thật nằm ở \textbf{Database Engine} (RBAC + VPD + OLS + Audit). Ứng dụng là tầng
phòng thủ bổ sung (defense-in-depth) + \textbf{mã hóa}: NNE bảo vệ \emph{đường truyền}, TDE bảo vệ
\emph{dữ liệu at-rest}. Tức là đồ án có đủ \textbf{cả Access Control lẫn Cryptography}.
\end{keybox}

\begin{center}
\resizebox{0.96\textwidth}{!}{%
\begin{tikzpicture}[font=\small,
  wide/.style={rectangle,rounded corners,draw,align=center,minimum width=13.4cm,minimum height=0.95cm},
  sub/.style={rectangle,rounded corners,draw=mainblue!60!black,fill=softblue,align=center,text width=2.7cm,minimum height=1.2cm,font=\footnotesize},
  arrow/.style={-{Stealth[length=3mm]},very thick}]
  \node[wide,fill=gray!10] (users) at (0,5.4) {\textbf{Người dùng}: ĐPV $\cdot$ Bác sĩ $\cdot$ KTV $\cdot$ Bệnh nhân $\cdot$ OLS (u1--u8) $\cdot$ DBA};
  \node[wide,fill=softblue] (app) at (0,3.8) {\textbf{Tầng ứng dụng (HospitalApp)}: InputValidator $\cdot$ mask CCCD $\cdot$ OracleErrorMapper $\cdot$ khóa đăng nhập $\cdot$ đổi mật khẩu};
  \node[rectangle,rounded corners,draw=red!60!black,fill=red!5,minimum width=13.4cm,minimum height=2.0cm] (db) at (0,1.15) {};
  \node[font=\bfseries,text=red!55!black] at (0,1.95) {Oracle Database Engine — ranh giới bảo mật thật};
  \node[sub] at (-4.9,0.75) {RBAC\\view + trigger\\(KTV, BN)};
  \node[sub] at (-1.63,0.75) {VPD\\lọc dòng\\(ĐPV, BS)};
  \node[sub] at (1.63,0.75) {OLS\\nhãn\\(THONGBAO)};
  \node[sub] at (4.9,0.75) {Audit + FGA\\ghi vết};
  \node[rectangle,rounded corners,draw=green!50!black,fill=green!8,align=center,text width=5.5cm,font=\footnotesize] (tde) at (0,-1.35) {\textbf{TDE} — mã hóa cột at-rest: CCCD, CMND, DIUNGTHUOC};
  \draw[arrow] (users)--(app);
  \draw[arrow] (app)--(db) node[midway,right,font=\footnotesize]{\;Oracle Net: NNE (AES256 + SHA512)};
  \draw[arrow,green!50!black] (db)--(tde) node[midway,right,font=\footnotesize]{\;ghi xuống đĩa};
\end{tikzpicture}}
\end{center}

\section{Phân công tổng quan}

\begin{center}
\resizebox{0.98\textwidth}{!}{%
\begin{tikzpicture}[font=\footnotesize, node distance=0.55cm,
  p/.style={rectangle,rounded corners,draw=mainblue!60!black,fill=softblue,align=center,text width=2.3cm,minimum height=1.5cm},
  arrow/.style={-{Stealth[length=2.4mm]},thick}]
  \node[p](a){\textbf{A}\\Mở đầu + Admin + TC\#1};
  \node[p,right=of a](b){\textbf{B}\\RBAC:\\KTV + BN};
  \node[p,right=of b](c){\textbf{C}\\VPD:\\ĐPV + BS};
  \node[p,right=of c](d){\textbf{D}\\OLS + Audit};
  \node[p,right=of d](e){\textbf{E}\\Backup/Recovery + Mã hóa};
  \draw[arrow](a)--(b); \draw[arrow](b)--(c); \draw[arrow](c)--(d); \draw[arrow](d)--(e);
\end{tikzpicture}}
\end{center}

\small
\begin{longtable}{c c L{6.6cm} L{3.7cm} c}
\toprule
\textbf{STT} & \textbf{Người} & \textbf{Phần phụ trách} & \textbf{Tài khoản chính} & \textbf{Thời lượng} \\
\midrule
\endhead
1 & \textbf{A} & Mở đầu + Kiến trúc + Phân hệ 1 (DBA) + \textbf{TC\#1} + Đổi mật khẩu & \code{SYSTEM} / \code{HOSPITAL_DBA} & 4 phút \\
2 & \textbf{B} & \textbf{YC1 -- RBAC}: Kỹ thuật viên + Bệnh nhân (view + trigger) & \code{KTV_NV006}, \code{BN_BN001} & 4 phút \\
3 & \textbf{C} & \textbf{YC1 -- VPD}: Điều phối viên + Bác sĩ & \code{DPV_NV001}, \code{BS_NV003} & 5 phút \\
4 & \textbf{D} & \textbf{YC2 -- OLS} + \textbf{YC3 -- Audit} & \code{u4_nvtk_hcm}, DBA & 5 phút \\
5 & \textbf{E} & \textbf{YC4 -- Backup/Recovery} + \textbf{Mã hóa (NNE + TDE)} + Kết luận & BVADMIN, DBA & 6 phút \\
\bottomrule
\end{longtable}
\normalsize

\begin{infobox}
\textbf{Tài khoản dùng chung:} mật khẩu nhân viên/bệnh nhân = \code{BV@2025!}; OLS \code{u4_nvtk_hcm} = \code{U4@2025};
DBA = \code{SYSTEM/oracle} hoặc \code{HOSPITAL_DBA/Hospital@DBA2025}; SYS = \code{<SYS_PASSWORD>} (tự điền).
\end{infobox}

\section{Chuẩn bị trước demo (checklist)}

\small
\begin{longtable}{L{13.2cm} c}
\toprule
\textbf{Việc cần chuẩn bị} & \textbf{Đã xong} \\
\midrule
\endhead
DB đang chạy, đã chạy migration 01--12 + \code{setup_all} (\code{run_migrations.ps1} / \code{scripts/setup.ps1}). & $\square$ \\
Đặt \code{NLS_LANG=.AL32UTF8} ở mọi cửa sổ PowerShell/sqlplus (tránh hỏng tiếng Việt). & $\square$ \\
Đã chạy \code{13_TDE_Encryption.sql} + keystore \textbf{OPEN} (kiểm \code{V$ENCRYPTION_WALLET}). & $\square$ \\
Đã bật \textbf{Native Network Encryption} trong \code{sqlnet.ora} (xem \file{SETUP_ENCRYPTION.md}). & $\square$ \\
Build app: \code{dotnet run --project HospitalApp} (đóng bản cũ trước). & $\square$ \\
Mở sẵn 2 cửa sổ: HospitalApp + SQL*Plus/SQL Developer (để chạy query DBA). & $\square$ \\
Dữ liệu mẫu nguyên vẹn (\code{SELECT COUNT(*) FROM BVADMIN.HSBA_DV;} = 5). & $\square$ \\
Đăng nhập thử trước 5 tài khoản để chắc không kẹt mật khẩu. & $\square$ \\
\bottomrule
\end{longtable}
\normalsize

\section{Người 1 (A) — Mở đầu + Kiến trúc + TC\#1 + Đổi mật khẩu}

\textbf{Nói:} giới thiệu bài toán (bệnh viện đa cơ sở, dữ liệu nhạy cảm), nêu \textbf{thông điệp xuyên suốt}.

\noindent\textbf{Demo:}
\begin{enumerate}
  \item \textbf{Phân hệ 1 -- AdminDashboard} (\code{SYSTEM/oracle} hoặc \code{HOSPITAL_DBA/Hospital@DBA2025}):
    tạo user demo \code{TEST_USER}, tạo role, grant/revoke role + object/system privilege; xem lại
    system / object / role privileges.
  \item \textbf{TC\#1 -- ánh xạ tài khoản:} mỗi người dùng app = một Oracle account riêng (cột
    \code{ORACLE_USER} trong \code{NHANVIEN}/\code{BENHNHAN}) \ra{} VPD/OLS tự áp theo \code{SESSION_USER}.
    Đăng nhập 2 vai trò khác nhau để thấy giao diện đổi theo role.
  \item \textbf{Đổi mật khẩu (self-service):} bấm nút \textbf{Đổi mật khẩu} ở header \ra{} nhập cũ/mới \ra{}
    ``Vui lòng đăng nhập lại'' \ra{} đăng nhập bằng mật khẩu mới.
\end{enumerate}

\begin{notebox}[title={Nói khi demo đổi mật khẩu}]
Dùng \code{ALTER USER ... IDENTIFIED BY ... REPLACE ...} — user tự đổi mật khẩu của chính mình
\textbf{không cần quyền DBA}. Sau khi đổi, connection cũ giữ mật khẩu cũ nên app đăng xuất để đăng nhập lại.
\end{notebox}

\begin{warnbox}[title={Có thể bị hỏi}]
\textit{``Brute-force tầng app khác gì DB profile?''} \ra{} App khóa sớm (5 lần sai/60s) trước khi gửi
kết nối lỗi xuống DB; DB profile là lớp chặn cuối.
\end{warnbox}

\section{Người 2 (B) — YC1 RBAC: Kỹ thuật viên + Bệnh nhân}

\textbf{Nói:} KTV và BN dùng \textbf{view + trigger} (RBAC) vì chỉ cần lọc theo 1 cột và giới hạn cột được sửa.

\noindent\textbf{Demo KTV} (\code{KTV_NV006/BV@2025!}):
\begin{enumerate}
  \item View \code{KTV_HSBA_DV_View} chỉ trả dịch vụ có \code{MAKTV = NV006} (không thấy của KTV khác).
  \item Cập nhật cột \textbf{KẾT QUẢ} \ra{} lưu thành công.
  \item (SQL*Plus) thử \code{UPDATE} cột khác (vd \code{LOAIDV}) \ra{} trigger chặn \textbf{\code{ORA-20001}}.
\end{enumerate}

\noindent\textbf{Demo Bệnh nhân} (\code{BN_BN001/BV@2025!}):
\begin{enumerate}
  \item Chỉ xem được thông tin của chính mình (view \code{BN_BENHNHAN_View} lọc theo \code{ORACLE_USER}).
  \item Sửa địa chỉ / tiền sử bệnh hợp lệ \ra{} lưu OK.
  \item Thử sửa \textbf{CCCD} \ra{} INSTEAD OF trigger chặn \textbf{\code{ORA-20002}} (không đổi được định danh).
  \item Chỉ ra \textbf{CCCD đang mask 4 số cuối} (\texttt{\textbullet\textbullet\textbullet\textbullet\textbullet\textbullet 5678}) — data minimization.
\end{enumerate}

\begin{warnbox}[title={Có thể bị hỏi}]
\textit{``Vì sao KTV dùng RBAC/view thay VPD?''} \ra{} KTV chỉ lọc theo \code{MAKTV}; view + trigger là đủ,
dễ kiểm soát cột update, dễ minh họa.
\end{warnbox}

\section{Người 3 (C) — YC1 VPD: Điều phối viên + Bác sĩ}

\textbf{Nói:} ĐPV/BS truy cập \textbf{cùng} các bảng \code{HSBA}/\code{HSBA_DV}/\code{DONTHUOC} nên dùng
\textbf{VPD} — Oracle tự chèn WHERE predicate vào mọi câu SQL, an toàn kể cả truy vấn trực tiếp.

\noindent\textbf{Demo Điều phối viên} (\code{DPV_NV001/BV@2025!}):
\begin{enumerate}
  \item \textbf{Thêm bệnh nhân mới} \ra{} app gọi \code{sp_create_benhnhan_full}: tự sinh MÃBN + tự tạo
    tài khoản \code{BN_<MABN>} (TC\#1). Báo lại mã + tài khoản đăng nhập.
  \item Tạo \textbf{HSBA} (mã sinh bằng sequence \code{fn_next_mahsba}), gán bác sĩ + gán KTV cho dịch vụ.
  \item Tab \textbf{Thông tin của tôi}: sửa quê quán / SĐT.
  \item Ô \textbf{CCCD}: mặc định mask, bấm \textbf{nút con mắt} để lộ đầy đủ khi đối chiếu (chỉ ĐPV mới có).
\end{enumerate}

\noindent\textbf{Demo Bác sĩ} (\code{BS_NV003/BV@2025!}):
\begin{enumerate}
  \item \textbf{VPD trên \code{HSBA}}: chỉ thấy hồ sơ có \code{MABS = NV003} (không thấy của BS khác).
  \item Cập nhật chẩn đoán/điều trị/kết luận \ra{} trigger ghi \code{LOG_BS_HSBA} (giá trị cũ/mới).
  \item Thêm dịch vụ chẩn đoán, chọn KTV; thêm/sửa đơn thuốc \ra{} \textbf{FGA + trigger} ghi vết.
  \item (tùy) cập nhật tiền sử/dị ứng của BN mình điều trị (TC\#3d) — column-grant.
\end{enumerate}

\begin{warnbox}[title={Có thể bị hỏi}]
\textit{``\code{update_check=>TRUE} để làm gì?''} \ra{} sau update, Oracle kiểm lại policy, chặn việc sửa
khóa-lọc để đẩy dữ liệu sang phạm vi khác. (Bài này để FALSE để tránh \code{ORA-28138} khi INSERT qua
INSTEAD OF — giải thích được nếu bị hỏi.)
\end{warnbox}

\section{Người 4 (D) — YC2 OLS + YC3 Audit}

\noindent\textbf{Demo OLS} (\code{u4_nvtk_hcm/U4@2025}):
\begin{enumerate}
  \item Mở \textbf{OLSViewerForm}, nhãn người dùng = \textbf{\code{NV:HCM:TK}}.
  \item Chỉ thấy thông báo có nhãn phù hợp, không thấy thông báo cấp cao hơn.
  \item (đối chứng) đăng nhập \code{u1_giamdoc} (cấp \textbf{BGD}) \ra{} thấy nhiều thông báo hơn.
  \item \textbf{Nói:} chính sách \code{BV_LABEL_POLICY} = \textbf{level} \code{NV < LDK < BGD} + \textbf{compartment}
    \code{HCM/HPN/HNI} + \textbf{group} \code{TH/TK/TM}.
\end{enumerate}

\noindent\textbf{Demo Audit} (DBA / SQL*Plus):
\begin{Code}
SELECT USERNAME, ACTION_NAME, OBJ_NAME, TIMESTAMP
FROM DBA_AUDIT_TRAIL ORDER BY TIMESTAMP DESC FETCH FIRST 20 ROWS ONLY;

SELECT DB_USER, OBJECT_NAME, POLICY_NAME, SQL_TEXT, EXTENDED_TIMESTAMP
FROM DBA_FGA_AUDIT_TRAIL ORDER BY EXTENDED_TIMESTAMP DESC FETCH FIRST 20 ROWS ONLY;
\end{Code}
\begin{itemize}
  \item Chỉ ra dòng audit ứng với thao tác Người 3 vừa làm (sửa chẩn đoán, thêm đơn thuốc).
  \item Mở bảng log nghiệp vụ: \code{LOG_BS_HSBA}, \code{LOG_BS_DONTHUOC}, \code{LOG_KTV_KETQUA}.
\end{itemize}

\begin{warnbox}[title={Có thể bị hỏi}]
\textit{``compartment vs group khác gì?''} \ra{} compartment nghĩa \textbf{AND} (phải có đủ); group nghĩa
\textbf{OR} và có \textbf{phân cấp}.
\end{warnbox}

\section{Người 5 (E) — YC4 Backup/Recovery + Mã hóa + Kết luận}

\subsection{YC4 -- Flashback recovery (chính, trực quan)}
\begin{Code}
$env:NLS_LANG = ".AL32UTF8"
sqlplus /nolog "@d:\repos\Oracle\PhanHe2\extras\recovery_demo.sql"
\end{Code}
Chỉ ra 3 mốc: \code{HSBA_DV} = \textbf{5} \ra{} xóa nhầm HS001 còn \textbf{4} \ra{} \textbf{FLASHBACK TABLE}
về \textbf{5} $\checkmark$. So sánh 3 phương pháp (xem \code{07_YC4_Backup_Recovery.sql} phần D): \textbf{RMAN}
(vật lý, toàn vẹn cao) $\cdot$ \textbf{Data Pump} (logic, di động) $\cdot$ \textbf{Flashback} (phục hồi cực nhanh
lỗi người dùng). RMAN/Data Pump cho xem lệnh + log; Flashback demo trực tiếp.

\subsection{Mã hóa đường truyền (NNE)}
\begin{Code}
SELECT NETWORK_SERVICE_BANNER
FROM   V$SESSION_CONNECT_INFO
WHERE  SID = SYS_CONTEXT('USERENV','SID');
\end{Code}
Chỉ ra dòng có \textbf{AES256} + \textbf{SHA} \ra{} kết nối TCP đã mã hóa (cấu hình \code{sqlnet.ora},
\code{ENCRYPTION_SERVER=REQUIRED}).

\subsection{Mã hóa cột at-rest (TDE)}
\begin{Code}
SELECT TABLE_NAME, COLUMN_NAME, ENCRYPTION_ALG
FROM   DBA_ENCRYPTED_COLUMNS ORDER BY TABLE_NAME, COLUMN_NAME;
\end{Code}
\code{CCCD}, \code{CMND} (NO SALT \ra{} giữ UNIQUE), \code{DIUNGTHUOC} (có salt) — \textbf{AES}. App vẫn đọc
plaintext (giải mã trong suốt); trên đĩa là ciphertext.

\begin{keybox}[title={Kết luận}]
Đồ án phủ \textbf{Access Control} (RBAC + VPD + OLS) + \textbf{Audit} + \textbf{Backup/Recovery} +
\textbf{Cryptography} (NNE đường truyền, TDE at-rest, masking) — bảo vệ dữ liệu y tế nhạy cảm nhiều tầng.
\end{keybox}

\begin{warnbox}[title={Có thể bị hỏi}]
\textit{``Vì sao cần Oracle Net encryption khi đã có RBAC/VPD?''} \ra{} RBAC/VPD/OLS bảo vệ \emph{bên trong}
database; NNE bảo vệ dữ liệu \emph{trên đường truyền} TCP. \quad \textit{``TDE có chặn ĐPV xem CCCD không?''}
\ra{} KHÔNG — TDE trong suốt với phiên hợp lệ; chặn xem là việc của access control/masking.
\end{warnbox}

\section{Luồng thời gian \& lưu ý}

Thứ tự: A (4') \ra{} B (4') \ra{} C (5') \ra{} D (5') \ra{} E (6'), tổng khoảng 24 phút + Q\&A. Mỗi người tự
kết nối tài khoản của mình; chuẩn bị sẵn cửa sổ để chuyển mượt.

\begin{riskbox}[title={Nếu thiếu thời gian}]
Rút gọn phần RMAN/Data Pump của Người 5 (chỉ nói, không chạy), giữ \textbf{Flashback + TDE + NNE} vì đây
là 3 minh chứng trực quan nhất cho phần mã hóa và phục hồi.
\end{riskbox}

\end{document}
